\documentclass[11pt]{article}
\usepackage[utf8]{inputenc}
\usepackage{amsmath,amssymb}
\usepackage[margin=1in]{geometry}
\usepackage{hyperref}
\emergencystretch=3em

\title{Headline statements, part III: the constant-Gaussian fluctuation model}
\author{Claude, with Daniel Murfet}
\date{2026-09-07}

\begin{document}
\maketitle
\sloppy

\begin{abstract}
Paper-facing wrappers for units~161 and 163 (\texttt{Laplace/Grammar/HeadlineGaussian.lean}):
\texttt{headline\_gaussian\_dichotomy} states that for $X\sim N(0,v)$ the expected leading moment
$\mathbb E_+[J_p(X)]$ is finite iff $\beta v<2$, equals $J_p(0)(1-\beta v/2)^{-p/2}$ below the
threshold and $+\infty$ above it; \texttt{headline\_gaussian\_jensen\_gap} states the strict
inequality $\mathbb E_+[J_p(X)]>J_p(0)$ for $v>0$. The module docstring fixes the scope (a Gaussian
constant fluctuation model of the leading coefficient's positive factor; nothing about the empirical
process or the fixed-$n$ chart integral). Zero \texttt{sorry}/\texttt{axiom}.
\end{abstract}

\section{Paper-facing meaning}
rem:pop\_vs\_emp~(ii) asserts $\mathbb E_{\mathcal D_n}[C(\xi_n)]\ne C(0)$ because the coefficients
depend nonlinearly on the fluctuation. The wrappers record the two precise consequences for the
simplest model: a Jensen gap and a temperature threshold $\beta v=2$ beyond which the expected
leading coefficient is infinite. The mirror \texttt{grammar\_lean.tex} already carries a sentence
to this effect in rem:pop\_vs\_emp; a dot is not placed since the remark is prose (bounded-statement
convention).

\section{Lean notes}
Wrappers only; compiled on the first attempt. Separate file since the Gaussian units import
\texttt{HeadlinePosterior.lean}.

\end{document}
